% Additive excerpt for the real-order chapter.
% Requires the manuscript's theorem/proof environments, amsmath and bibliography.
% Full proofs and exact certificates: the accompanying research article.
\section{Two proved global normalized-radius regimes}
\label{grm:sec:global-radius}

The local coefficient theorem does not exclude a later change of radial
sign. The following two full-radius statements give explicit regions
where no such change occurs; their complete analytic proofs and finite
rational certificates are in \cite{grm:report}.

\begin{theorem}[Uniform outer-order region]
\label{grm:thm:outer}
For every real $a\ge8$, every $b>0$, and $0<\rho\le1$,
\[
 \frac{d}{d\rho}\left(\frac{\cos\theta_{a,b}(\rho)}\rho\right)
 <-\frac{\rho H_4^{(b)}}{120}\left(\frac25\right)^a<0.
\]
Moreover,
\[
 \frac{1+2^{-b}}2\left(\frac23\right)^a
 -\frac{\cos\theta_{a,b}(\rho)}\rho
 >\frac{\rho^2 H_4^{(b)}}{240}\left(\frac25\right)^a.
\]
\end{theorem}
\begin{proof}[Proof summary; complete proof in the cited report]
Set $s=\rho^2$, $m=2\cos\theta/\rho$, and
$c_n=H_{n-1}^{(b)}/n^a$. The zero equation is
$P(m,s)=\sum_{n\ge2}c_nT_n(m,s)=0$, with
$T_2=m$, $T_3=sm^2-1$, and $T_{n+1}=s(mT_n-T_{n-1})$.
Uniform coefficient and Chebyshev estimates place every root in
$0<m<(11/10)c_3/c_2<1/10$ and give
$(4/5)c_2<P_m<(6/5)c_2$.
A degree-four Bernstein certificate gives
$\sum_{n=2}^6c_n\partial_sT_n>(2/5)c_5$; an analytic infinite-tail
bound gives $|\sum_{n\ge7}c_n\partial_sT_n|<(39/100)c_5$.
Hence $m'=-P_s/P_m<-c_5/(120c_2)$.
Since $(\cos\theta/\rho)'=\rho m'$, the first bound follows;
integration gives the second. The inequalities extend from $a=8$ to
all real $a\ge8$ by explicit exponential comparisons, not sampling.
\end{proof}

\begin{theorem}[Explicit finite-inner-order transfer]
\label{grm:thm:inner}
Let $a>0$ and set
$\Delta(a)=(2/5)^a-(3/8)^a$.
If $2^{-b}<\Delta(a)/256$, then $\cos\theta_{a,b}(\rho)/\rho$
strictly decreases on $0<\rho\le1$.
If $2^{-b}\le\Delta(a)/512$, its derivative is at most
$-\rho\Delta(a)/16$.
\end{theorem}
\begin{proof}[Proof summary; complete proof in the cited report]
Write the positive double kernel using independent variables
$X=e^{-Y}$, $V=e^{-Z}$, with $Y$ gamma of shape $a$, rate 2, and $Z$
gamma of shape $b$, rate 1. The zero equation, up to a positive factor,
is $G_b=\mathbb E[(m-X-XV)/(D_XD_{XV})]=0$, where
$D_u=1+s(u^2-mu)$.
On $m,s\in[0,1]$, the equation and its two first derivatives differ
from the $V=0$ equation by at most $12\,2^{-b}$, $64\,2^{-b}$, and
$64\,2^{-b}$, respectively. The limiting derivatives obey
$(G_\infty)_m\ge1/4$ and
$(G_\infty)_s\ge[\mathbb EX^3-(\mathbb EX^2)^2/\mathbb EX]/4
=\Delta(a)/4$.
Endpoint signs and the inherited angular uniqueness theorem identify
the root; the stated conditions make both derivatives positive there.
Implicit differentiation proves the claims.
\end{proof}

For integer $a=1,\ldots,7$, sufficient integer lower bounds on $b$
are respectively $14,14,15,16,17,18,19$.
Thus the integer conjecture for $a\ge2$ remains open here only on six
bounded inner-order strips with $a=2,\ldots,7$. This is not a claim
that those remaining strips have been certified.

\subsection{A positive binomial identity}
\label{grm:sec:binomial}
The positive double kernel can be averaged over the interval between
its two nodes, giving a positive measure $\mu_{a,b}$ such that
\[
 F_{a,b}(z)=z^2\int_0^1\frac{d\mu_{a,b}(u)}{(1-uz)^2},
 \qquad \int u^k\,d\mu_{a,b}(u)=\frac{c_{k+2}}{k+1}.
\]
Consequently
\[
 \mathcal B_k(a,b)=\sum_{j=0}^k(-1)^j\binom kj
       \frac{H_{2j+2}^{(b)}}{(2j+2)(2j+3)^a}
 =\int u(1-u^2)^k\,d\mu_{a,b}(u)>0.
\]
For all $a,b>0$,
\[
 -\operatorname{Im}F_{a,b}(i)
 =\sum_{k\ge0}\frac{k+1}{2^{k+1}}\mathcal B_k(a,b).
\]
After $N$ terms the positive remainder is strictly less than
$(N+2)H_2^{(b)}/(2^{N+1}3^a)$.
This is a distinct positive transform, not a replacement claim for the
sharp Euler evaluator or an arithmetic $S_6$/$S_8$ reduction.
